\chapter{Regular 5-Story Building Example}
\label{ch:regular-5story-building}

\section{Purpose}

This chapter documents the \emph{elastic baseline} mixed-building example that
preceded the materially nonlinear reinforced-concrete continuation developed in
Chapter~\ref{ch:regular-5story-rc-nonlinear-building}.  It remains useful as
the simplest reference configuration for understanding the shared
\texttt{Domain}, mixed 1D/2D topology, and polymorphic structural wrapping
without the additional complexity of path-dependent constitutive response.

The example is a regular reinforced-concrete-like building with:
\begin{itemize}
  \item five stories,
  \item four grid axes in the \(x\)-direction,
  \item three grid axes in the \(y\)-direction,
  \item frame members modelled with 3D Timoshenko beams,
  \item slabs modelled with MITC4 / Mindlin--Reissner shells.
\end{itemize}

The objective is not a code-check design, but a pedagogical and architectural
integration example:
\begin{itemize}
  \item shared geometry and topology in one \texttt{Domain},
  \item mixed 1D/2D structural families,
  \item one polymorphic structural solver container,
  \item one structural \texttt{.vtm} export for ParaView.
\end{itemize}

\section{Unit System}

The example is written consistently in:
\[
\text{force} = \mathrm{kN},
\qquad
\text{length} = \mathrm{m},
\qquad
\text{stress} = \mathrm{kN/m^2}.
\]

Therefore:
\begin{itemize}
  \item Young's moduli are entered in \(\mathrm{kN/m^2}\),
  \item slab loads are entered in \(\mathrm{kN/m^2}\),
  \item nodal lateral forces are entered in \(\mathrm{kN}\),
  \item computed displacements are reported in \(\mathrm{m}\).
\end{itemize}

\section{Geometry}

\subsection{Plan and height}

The plan grid is:
\[
x = \{0,\; 6,\; 12,\; 18\}\ \mathrm{m},
\qquad
y = \{0,\; 5,\; 10\}\ \mathrm{m}.
\]

This gives:
\begin{itemize}
  \item \(3\) bays of \(6\,\mathrm{m}\) in the \(x\)-direction,
  \item \(2\) bays of \(5\,\mathrm{m}\) in the \(y\)-direction.
\end{itemize}

The storey height is:
\[
h = 3.20\ \mathrm{m},
\]
so the floor elevations are:
\[
z_k = 3.20\,k, \qquad k = 0,\dots,5.
\]

The total structural height is:
\[
H = 16.0\ \mathrm{m}.
\]

\subsection{Discretisation}

All members share the same building nodes:
\begin{itemize}
  \item \(4 \times 3 = 12\) plan nodes per level,
  \item \(6\) levels including the base,
  \item total node count:
        \[
        n_\mathrm{nodes} = 12 \times 6 = 72.
        \]
\end{itemize}

The element counts are:
\begin{itemize}
  \item columns:
        \[
        n_c = 5 \times 4 \times 3 = 60,
        \]
  \item beams:
        \[
        n_b = 5 \left[(4-1) \times 3 + (3-1) \times 4\right] = 85,
        \]
  \item slab panels:
        \[
        n_s = 5 \times (4-1) \times (3-1) = 30.
        \]
\end{itemize}

Hence the total number of structural geometric elements is:
\[
n_\mathrm{elem} = 60 + 85 + 30 = 175.
\]

\section{Sections and Materials}

\subsection{Columns}

Columns are modelled as 3D Timoshenko beams with rectangular section:
\[
b_c = 0.45\ \mathrm{m},
\qquad
h_c = 0.45\ \mathrm{m}.
\]

The frame material is idealised as linear elastic with:
\[
E_f = 28.0 \times 10^6\ \mathrm{kN/m^2},
\qquad
\nu = 0.20.
\]

The shear modulus is reconstructed as:
\[
G_f = \frac{E_f}{2(1+\nu)}.
\]

\subsection{Beams}

Beams are also modelled as 3D Timoshenko beams, with section:
\[
b_b = 0.30\ \mathrm{m},
\qquad
h_b = 0.60\ \mathrm{m}.
\]

The same frame elastic modulus is used:
\[
E_b = E_f = 28.0 \times 10^6\ \mathrm{kN/m^2}.
\]

For both columns and beams the section properties are computed directly from
the rectangular dimensions:
\[
A = b h,
\qquad
I_y = \frac{b h^3}{12},
\qquad
I_z = \frac{h b^3}{12}.
\]

The torsional constant is approximated with the standard rectangular
Saint-Venant estimate:
\[
J \approx \frac{b_\text{min}^3 h_\text{max}}{3}
\left(1 - 0.63 \frac{b_\text{min}}{h_\text{max}}\right).
\]

The Timoshenko shear factors are taken as:
\[
\kappa_y = \kappa_z = \frac{5}{6}.
\]

\subsection{Slabs}

Each floor slab panel is modelled by a MITC4 / Mindlin--Reissner shell with
thickness:
\[
t = 0.15\ \mathrm{m}.
\]

The slab elastic modulus is slightly reduced with respect to the frame:
\[
E_s = 25.0 \times 10^6\ \mathrm{kN/m^2},
\qquad
\nu_s = 0.20.
\]

\section{Loads and Boundary Conditions}

\subsection{Boundary conditions}

All base nodes, that is all nodes on the plane \(z=0\), are fully fixed.
This is implemented at the model level through:
\begin{verbatim}
model.fix_z(0.0);
\end{verbatim}
which, in the current API, constrains all DOFs of every node lying on that
plane.

\subsection{Gravity loading}

Each slab receives a uniform downward pressure:
\[
q_g = -7.50\ \mathrm{kN/m^2}.
\]

This pressure represents a simplified vertical service load
(self-weight + superimposed permanent load + a modest live-load component).

The load is assembled consistently from the shell interpolation:
\[
\mathbf{f}_I = \int_{\Omega_e} N_I\,\mathbf{t}\,dA,
\qquad
\mathbf{t} = (0,0,q_g)^T.
\]

\subsection{Horizontal loading}

To make frame bending visually more evident, one corner node of every floor
receives a horizontal force in the global \(x\)-direction. The loaded node at
storey \(k\) is:
\[
(x,y,z) = (18,\; 10,\; z_k).
\]

The storey force pattern is triangular:
\[
F_x^{(k)} = 400\,k \quad \mathrm{kN},
\qquad
k = 1,\dots,5.
\]

Thus the applied forces are:
\[
400,\ 800,\ 1200,\ 1600,\ 2000\ \mathrm{kN}.
\]

This is intentionally not a perfectly symmetric wind model. The asymmetry is
useful here because it activates:
\begin{itemize}
  \item global building sway,
  \item perimeter beam bending,
  \item slab diaphragm action,
  \item a modest torsional component.
\end{itemize}

\section{Architectural Realisation in Code}

\subsection{One shared \texttt{Domain}}

The first important architectural point is that the example inserts
\emph{all} structural geometry into the same \texttt{Domain<3>}:
\begin{itemize}
  \item columns are stored as 1D line geometries,
  \item beams are stored as 1D line geometries,
  \item slabs are stored as 2D quadrilateral geometries.
\end{itemize}

This is exactly the kind of mixed topology that motivated the recent
\texttt{DMPlex} hardening described in Chapter~\ref{ch:dmplex-mixed-topology}.
The resulting \texttt{DMPlex} remains non-interpolated, but it is now
semantically correct for mixed 1D/2D structural families.

For this model the reported \texttt{DMPlex} dimension is:
\[
\dim(\texttt{DM}) = 2,
\]
because the maximum topological dimension in the domain is the slab topology.

\subsection{Polymorphic structural container}

The solver model is assembled as:
\[
\texttt{Model<TimoshenkoBeam3D,\ SmallStrain,\ 6,\ SingleElementPolicy<StructuralElement>>}.
\]

This means:
\begin{itemize}
  \item the element container is
        \[
        \texttt{std::vector<StructuralElement>},
        \]
  \item each concrete geometry is wrapped into the appropriate concrete
        element type:
        \begin{itemize}
          \item \texttt{BeamElement<TimoshenkoBeam3D,3>} for columns,
          \item \texttt{BeamElement<TimoshenkoBeam3D,3>} for beams,
          \item \texttt{ShellElement<MindlinReissnerShell3D>} for slabs;
        \end{itemize}
  \item the global assembly sees a homogeneous interface, while the structural
        VTK exporter can still dispatch back to the concrete formulation
        through \texttt{StructuralReductionPolicy<ElementT>}.
\end{itemize}

\subsection{Why this example is a good architectural integrator}

This one example exercises, simultaneously:
\begin{itemize}
  \item geometric ownership in \texttt{Domain},
  \item mixed-topology \texttt{DMPlex} assembly,
  \item analysis-node DOF assignment on shared vertices,
  \item heterogeneous structural element wrapping,
  \item shell consistent-load assembly,
  \item formulation-aware structural reconstruction,
  \item multi-block VTK output.
\end{itemize}

In that sense it is a better ``main example'' than the previous collection of
independent benchmarks: it shows how the pieces work together.

\section{Execution Results}

The executed example stored in the repository reported:
\[
\max |u| = 9.248 \times 10^{-3}\ \mathrm{m},
\]
\[
\max |u_x| = 8.292 \times 10^{-3}\ \mathrm{m},
\qquad
\max |u_z| = 2.523 \times 10^{-3}\ \mathrm{m}.
\]

At the loaded roof corner node, the displacement vector was:
\[
u_\mathrm{roof} =
\begin{bmatrix}
  8.292 \times 10^{-3} \\
 -3.598 \times 10^{-3} \\
 -1.955 \times 10^{-3}
\end{bmatrix}
\mathrm{m}.
\]

The roof drift ratio in the \(x\)-direction is therefore:
\[
\frac{u_x^\mathrm{roof}}{H}
=
\frac{8.292 \times 10^{-3}}{16.0}
= 5.18 \times 10^{-4}.
\]

These values are small enough to remain in the linear range, but large enough
to make the deformed shape and beam flexure visible in ParaView when a moderate
\texttt{Warp By Vector} scale factor is used.

\section{VTK Output}

The example writes the following VTK files:
\begin{itemize}
  \item \texttt{regular\_5story\_building\_structural.vtm},
  \item \texttt{regular\_5story\_building\_structural\_axis.vtp},
  \item \texttt{regular\_5story\_building\_structural\_surface.vtp},
  \item \texttt{regular\_5story\_building\_structural\_sections.vtp},
  \item \texttt{regular\_5story\_building\_structural\_material\_sites.vtp},
  \item \texttt{regular\_5story\_building\_structural\_fibers.vtp}.
\end{itemize}

The semantic meaning of the blocks remains the same as in the previous
structural exporter chapters:
\begin{description}
  \item[\texttt{axis}] line carriers for beam centre-lines and shell edges;
  \item[\texttt{surface}] reconstructed beam skins and shell faces;
  \item[\texttt{sections}] section stations and through-thickness carriers;
  \item[\texttt{material\_sites}] constitutive sampling sites;
  \item[\texttt{fibers}] fiber-section carriers (empty in this example).
\end{description}

In ParaView, the intended workflow is:
\begin{enumerate}
  \item open the \texttt{.vtm} file,
  \item display the \texttt{surface} block,
  \item apply \texttt{Warp By Vector} using \texttt{displacement},
  \item colour by reconstructed resultants or stresses,
  \item inspect \texttt{sections} and \texttt{material\_sites} for the
        formulation-specific internal fields.
\end{enumerate}

\section{Why the Example Matters}

From a software-architecture perspective, this example confirms that the
current direction is viable:
\begin{itemize}
  \item mixed structural families can share one geometric and topological
        domain,
  \item \texttt{StructuralElement} is sufficient as the polymorphic solver
        boundary,
  \item reconstruction remains formulation-specific and therefore reusable for
        later scale-transfer work,
  \item the new \texttt{main.cpp} now acts as a readable cookbook for building
        a realistic mixed frame--shell model.
\end{itemize}

That is exactly the kind of example needed before moving on to the next layer
of abstraction, namely a dedicated \texttt{DMPlex} policy or builder strategy.
